% !TeX root = ../../tesis.tex

% =============================================================================
%  B.6 — API REST: ENDPOINTS FASTAPI
% =============================================================================
\section{\api{} REST de VFTModel: Endpoints Disponibles}
\label{anx:vft-api}

El Cuadro~\ref{tab:vft-endpoints} lista los ocho endpoints de VFTModel,
agrupados por su función. Los de la capa \textit{network} calculan indicadores
topológicos escalares o tabulares; los de la capa \textit{geolayers} (§
\ref{anx:vft-geolayers}) serializan los mismos resultados como
\texttt{FeatureCollection} consumibles directamente por Transport-\textsc{gis}.

\begin{table}[H]
    \centering
    \small
    \caption[\api{} REST de VFTModel: endpoints disponibles]{Endpoints de la
    \api{} \textsc{rest} de VFTModel y el indicador o acción que producen.
    \fuentefigura{\texttt{src/api/main.py} y \texttt{src/api/routes/geo\_layers.py}
    (VFTModel, 2026).}}
    \label{tab:vft-endpoints}
    \begin{tabularx}{\textwidth}{|l|l|X|}
        \hline
        \textbf{Endpoint} & \textbf{Método} & \textbf{Resultado} \\
        \hline
        \multicolumn{3}{|l|}{\textit{Capa topológica (\texttt{/api/v1/network/})}} \\
        \hline
        \texttt{/build-auto}
            & GET  & Construye el grafo y lo almacena en caché \\
        \texttt{/spatial-coverage}
            & GET  & Cobertura (\%) por alcaldía o municipio \\
        \texttt{/topological/capillary-strength}
            & GET  & Ranking de fuerza capilar por nodo \\
        \texttt{/topological/geo-capillary}
            & GET  & Detección de macro-\termino{hubs} por proximidad espacial \\
        \texttt{/topological/detour-factor}
            & GET  & Factor de desviación general o por par O-D \\
        \texttt{/topological/average-travel-time}
            & GET  & Tiempo de viaje promedio ($T$) \\
        \texttt{/topological/betweenness-centrality}
            & GET  & Centralidad de intermediación ($B(v)$) \\
        \hline
        \multicolumn{3}{|l|}{\textit{Capa GeoLayers (\texttt{/api/v1/network/geolayers/})}} \\
        \hline
        \texttt{/coverage}
            & GET  & Cobertura como \texttt{FeatureCollection} GeoJSON \\
        \texttt{/capilar}
            & GET  & Fuerza capilar con geometría de nodos \\
        \texttt{/detour}
            & GET  & Factor de desviación con trayectorias O-D \\
        \hline
    \end{tabularx}
\end{table}

\subsection*{Parámetros del endpoint \texttt{build-auto}}

Este endpoint es el punto de entrada obligatorio antes de calcular cualquier
indicador: construye el grafo, aplica la impedancia y almacena el resultado en
un caché en memoria (\texttt{Singleton}). Los parámetros configurables son el
modo de construcción y la tolerancia de transbordo:

\begin{lstlisting}[style=estiloPython,
    caption={Endpoint de construcción y calentamiento del caché del grafo.
             \texttt{src/api/main.py}.},
    label=lst:vft-build-auto]
@app.get("/api/v1/network/build-auto")
async def build_network_auto(
    mode:        str   = Query("REALISTIC_INTEGRATION",
                               description="STRICT_TOPOLOGY | REALISTIC_INTEGRATION"),
    tolerance_m: float = Query(85.0,
                               description="Tolerancia peatonal de transbordo en metros")
):
    G = await get_or_build_graph(mode, tolerance_m)
    return {
        "status":  "success",
        "nodos":   G.number_of_nodes(),
        "aristas": G.number_of_edges()
    }
\end{lstlisting}
